\subsection{\rev{repoDB normalization and cross-resource entity alignment}}
\label{sec:repodb_alignment}
\begin{revisionblock}
\textbf{Normalization and target definition.}
We removed HTML and excess whitespace, converted identifiers to uppercase, and extracted DrugBank IDs (\texttt{DB} followed by five digits) and UMLS CUIs (\texttt{C} followed by seven digits) from structured fields, using raw text fields when necessary. All 10,800 input rows contained valid identifiers. Drugs and diseases were represented as \texttt{Compound::DBxxxxx} and \texttt{Disease::UMLS:Cxxxxxxx}; names were used for display, not cross-resource matching. Approved records defined positive targets, while terminated, withdrawn, suspended, and failed records formed a separate diagnostic relation. Deduplication yielded 6,677 approved and 3,281 diagnostic drug--disease pairs. The 21 overlapping pairs retained their approved label, leaving 3,260 diagnostic pairs. These diagnostic labels were not treated as positive indications or mechanistic contraindications.

\textbf{Entity alignment.}
We linked compounds by exact canonical DrugBank IDs to the prepared DRKG entity map of 13,918 entities. Across approved and diagnostic pairs, 1,130 of 1,571 compounds (71.93\%) matched; among approved compounds, 1,106 of 1,519 matched (72.81\%). The 2,051 disease CUIs had no exact identifier matches and remained local UMLS nodes. No ontology crosswalk, synonym-based disease alignment, or disease-to-gene conversion was applied. Unmatched compounds were also retained. Thus, local graph-ID coverage includes entities without imported DRKG evidence, while the alignment percentages refer only to the prepared DRKG map, not the complete raw resource.

\textbf{Evidence graph construction.}
We combined training approved-indication edges and diagnostic edges with DRKG compound--gene edges linked through matched compounds, preserving relation types, directions, and gene IDs. Gene--gene edges were included only when both endpoints belonged to this gene set; no DRKG disease--gene edges were imported through a UMLS mapping. After deduplication, edges were capped in input order: 80 diagnostic edges per disease and per compound, 80 compound--gene edges per compound, and 40 gene--gene edges per gene with a total cap of 150,000. The resulting graph contained 5,677 approved, 3,222 diagnostic, 39,017 compound--gene, and 42,888 gene--gene edges: 90,804 triples over 8,818 entities and 63 relations. All task entities remained in the local ID maps. The reduction from 3,260 to 3,222 diagnostic pairs resulted from the edge caps.

\textbf{Splitting and leakage control.}
The 6,677 approved triples were divided into 5,677/500/500 training/validation/test triples using seed 2028. A triple was eligible for holdout only if both its compound and disease retained another approved training edge. The 1,519 training compounds defined the candidate universe. The splits were disjoint, all held-out entities remained represented in training, and held-out approved triples and their diagnostic counterparts were excluded from the evidence graph. These constraints prevent exact target-triple leakage but do not remove all possible semantic overlap between resources. Evaluation lists were generated without gold injection, with canonical and numeric IDs stored separately from display names. This preparation supplies drug-side molecular evidence and gene bridges while retaining a partially aligned disease context.
\end{revisionblock}
